%% ma: L10-B15-0098
%% lop: 10
%% chuong: 6
%% bai: 15
%% dang: 10.15.5
%% loai: TN4
%% muc: VD
%% dap_an: "A"
%% nguon: "Ảnh giáo viên gửi 25/11/2026 (lô 2): Bài tập trắc nghiệm – Toán 10 Bài 15 Hàm số (sự biến thiên, đồ thị), Phần I Câu 9"
%% kien_thuc: "Lớp 10 Bài 15 (đồng biến, nghịch biến); hàm số bậc nhất y=ax+b đồng biến khi a>0 (lớp 9)"
%% trang_thai: chinh-thuc
%% ghi_chu: "Biến thể thay số của L10-B15-0048. Phương án A gốc là 4040. Nếu muốn giữ 4040 thì đổi đoạn thành [-2022;2022]."
\begin{ex}
Có bao nhiêu giá trị nguyên của tham số $m$ thuộc đoạn $[-2022;2023]$ để hàm số $y=(m^2-4)x+2m$ đồng biến trên $\mathbb R$?
\choice{\True $4041$}{$4044$}{Vô số}{$2020$}
\loigiai{Cần $m^2-4>0\Leftrightarrow m<-2$ hoặc $m>2$. Trên $[-2022;2023]$: $m\in\{-2022;\dots;-3\}$ có $2020$ giá trị, $m\in\{3;\dots;2023\}$ có $2021$ giá trị. Tổng $4041$.}
\end{ex}
